\section{Context of $e^+e^-$ Alliance}
\label{sec:context}

Achieving Standard Model predictions at the level of precision required for future collider programs, including the High-Luminosity Large Hadron Collider (HL-LHC), is essential to ensure robust control of systematic uncertainties in the interpretation of potential new physics anomalies. The demand for such precision is driven by percent-level experimental uncertainties and the unprecedented statistical power of modern collider datasets. In this regime, theoretical uncertainties become limiting, and the dominant contribution arises from our modeling of Quantum Chromodynamics, often referred to as precision QCD.

Precision QCD studies can be pursued at hadron colliders such as the LHC, but they are most cleanly and systematically tested in lepton collider environments. At present, however, there is no active high-energy lepton collider operating at LEP energies or above. To address this, the $e^+e^-$ alliance~\cite{EEAlliance:website} has successfully curated and reanalyzed archived datasets from the Large Electron--Positron collider (LEP), the precursor to the LHC, opening new opportunities enabled by three decades of theoretical advances and modern experimental techniques. The combination of ultra-precise $e^+e^-$ data and state-of-the-art perturbative and non-perturbative QCD tools provides a uniquely powerful environment for high-precision studies. The alliance has produced a series of precision measurements using these archived datasets, including two-particle correlations~\cite{Badea:2019vey}, jet energy spectra and substructure~\cite{Chen:2021uws}, long-range near-side correlations~\cite{Chen:2024lrc}, energy-energy correlators~\cite{Bossi:2025eec}, and unbinned measurements of thrust~\cite{Badea:2025ubt}.

More broadly, this program is a compelling opportunity for AI for Science in a domain where precision is paramount. While AI is increasingly used to assist theoretical calculations and generate predictions, in the physical sciences such predictions ultimately require experimental validation. In high-energy physics, this validation depends on collider datasets that are rarely publicly accessible. The $e^+e^-$ environment, particularly at LEP energies, provides the cleanest setting for precision tests of QED and QCD, with well-defined initial states and calculable dynamics. The successful curation and use of archived open data therefore establishes the conditions for a theory--experiment loop in which AI systems, guided by domain experts, can propose new calculations, observables, or improved models and directly confront them with high-precision data. This creates a concrete framework for developing AI systems that do not merely interpolate existing knowledge, but actively refine theoretical descriptions through direct comparison with experiment.

\subsection*{Lessons from AI--Physicist Collaboration}

Carrying out this analysis has also produced practical insights into how AI agents and physicists can work together effectively. Our workflow followed a constrained human-in-the-loop cycle in which the physicist defined physics objectives, acceptance criteria, and interpretation authority, while the AI agent owned implementation throughput, consistency checks, and traceability. Crucially, experienced physicists also contributed domain intuition that cannot be reduced to formal specifications alone, including judgments about physically meaningful binning, acceptable regularization behavior, uncertainty communication, and when apparent anomalies warranted deeper investigation. Requests were made concrete and executable, each iteration linked to specific scripts, outputs, and note sections, and corrections were applied incrementally and revalidated rather than batched into large unreviewed changes.

From this experience, several recurring challenge classes emerged. Implicit physics priors, expert expectations about acceptable binning, regularization behavior, or uncertainty presentation, were often assumed rather than explicitly stated, requiring iterative clarification. Precision formatting constraints (label offsets, legend placement, pad geometry) required many short feedback loops, and apparent cosmetic changes sometimes affected interpretability and therefore physics conclusions. State synchronization across rapidly evolving scripts, outputs, and note text demanded continuous provenance refreshing. Throughout, authority remained intentionally asymmetric: the physicist owns physics choices and release signoff, while the agent executes and refactors. These lessons motivate stricter templating for figures, stronger configuration contracts, and clearer channels for encoding expert heuristics into machine readable constraints in future iterations of this program.
